\documentclass[11pt]{article}
\usepackage[margin=1in]{geometry}
\usepackage[T1]{fontenc}
\usepackage[utf8]{inputenc}
\usepackage{lmodern}
\usepackage{amsmath}
\usepackage{enumitem}
\usepackage{amsfonts}
\usepackage{amssymb}

\setlist[itemize]{noitemsep, topsep=2pt}
\setlist[enumerate]{noitemsep, topsep=2pt}

\title{TickRange Accrual Notes Replication}
\author{}
\date{\today}

\begin{document}
\maketitle
\tableofcontents


\section{Swaps}


\subsection{Occupation Time}
impliedOccupation time calculated from implied volatility

realizedOccupationTime


CFMM(RANGE_ACCRUAL_NOTE):
	tradingFunction(payoff(RANGE_ACCRUAL_NOTE)):
	rule: NoArbitrage: price(RANGE_ACCRUAL_NOTE) =  price(liquidityPosition(underlying))
        - price(futuresPerpetual(underlying))
						       	     + price(Option(underlying)

		rule NoArbitrage: price(RANGE_ACCRUAL_NOTE) =
	     		  sum(discounted(price(Option(underlying),tickLower))
			      - discounted(price(Option(underlying), tickUpper)))

		invariant (tradingFunction):



\subsection{Applications}

Consider a liquidity provider who wishes to hedge against volatility for his current liquidity position. With range structured notes, we can do that  the following way:

The LP is    feeGrowthInside -      feeGrowthOutside
                 [il, iu]        (i_min, i_u) U (i_u, i_max)              


                 The LP has access to  hitorical trading volume data and is currently located on the optimal range that captures most of the volume: [i_l*,i*_u] = [il,iu]. The LP is long theta, However:

                 - A tail risk often not consider and priced is the risk induced by how the volatility of the underlying shapes expected time inside the position. This is a risk that can be captured with TAN and thus, aa TAN can be used to hedge such risk

                 - Another risk is the PLP adversion to JIT liquidity. We show that a TAN can  be used by PLP to hedge against the JIT rent-extraction risk induced by


                 
\subsubsection{HEDGING JIT}
- Micrstructure context:

\partial cosOfLiquidity / \partial feeRevenueAllocatioRule 

Iff feeRevenueAllocationRule := \propto unitOfLiquidity (proRata AllocationRule) ==> 
                                                                 |
                                                                 |
                                                             (competition)
                                                                 |
                                                       (larger share of feeRevenue/tick)
                                                                 |
                      larger unitOfLiquidity/tick          <-------   
                                 |
                                 |
                                 v
                         (excess liquidity)
                      
\section{References}

- replicating market makers
- replicatios monotonic payoffs without oracles

\end{document}
